\appendix

\chapter{Cat\'alogo de fun\c{c}\~oes}
\label{ap:funcoes}

Esta se\c{c}\~ao lista as 43 fun\c{c}\~oes do c\'odigo Rust, agrupadas por
m\'odulo. O \emph{stub} \code{\_start} n\~ao \'e uma fun\c{c}\~ao Rust e
aparece separadamente no Cap\'itulo~\ref{cap:boot}.

\section{\code{main.rs} (2 fun\c{c}\~oes)}

\begin{longtable}{@{}p{6.2cm}p{8.3cm}@{}}
\caption{Fun\c{c}\~oes de \code{main.rs}.}\\
\toprule
\textbf{Assinatura} & \textbf{Descri\c{c}\~ao} \\
\midrule
\endfirsthead
\toprule
\textbf{Assinatura} & \textbf{Descri\c{c}\~ao} \\
\midrule
\endhead
\code{kernel\_main(multiboot\_magic: u32, multiboot\_info: u32) -> !} &
Ponto de entrada Rust; inicializa drivers, imprime banner e entra no shell. \\
\code{panic(info: \&PanicInfo) -> !} &
\emph{Panic handler}; imprime a mensagem e para a CPU. \\
\bottomrule
\end{longtable}

\section{\code{arch.rs} (4 fun\c{c}\~oes)}

\begin{longtable}{@{}p{5.2cm}p{9.3cm}@{}}
\caption{Fun\c{c}\~oes de \code{arch.rs}.}\\
\toprule
\textbf{Assinatura} & \textbf{Descri\c{c}\~ao} \\
\midrule
\endfirsthead
\toprule
\textbf{Assinatura} & \textbf{Descri\c{c}\~ao} \\
\midrule
\endhead
\code{inb(port: u16) -> u8} &
L\^e um byte de uma porta. \code{unsafe}. \\
\code{outb(port: u16, value: u8)} &
Escreve um byte em uma porta. \code{unsafe}. \\
\code{halt\_loop() -> !} &
La\c{c}o infinito com \code{cli; hlt}. \\
\code{reboot() -> !} &
Envia \code{0xFE} ao 8042; se falhar, para a CPU. \\
\bottomrule
\end{longtable}

\section{\code{sync.rs} (5 fun\c{c}\~oes)}

\begin{longtable}{@{}p{5.7cm}p{8.8cm}@{}}
\caption{Fun\c{c}\~oes de \code{sync.rs}.}\\
\toprule
\textbf{Assinatura} & \textbf{Descri\c{c}\~ao} \\
\midrule
\endfirsthead
\toprule
\textbf{Assinatura} & \textbf{Descri\c{c}\~ao} \\
\midrule
\endhead
\code{SpinLock::new(data: T) -> SpinLock<T>} &
Cria o lock em contexto \code{const}. \\
\code{SpinLock::lock(\&self) -> SpinLockGuard<'\_, T>} &
Adquire o lock girando at\'e conseguir. \\
\code{SpinLockGuard::deref(\&self) -> \&T} &
Acesso imut\'avel ao dado protegido. \\
\code{SpinLockGuard::deref\_mut(\&mut self) -> \&mut T} &
Acesso mut\'avel ao dado protegido. \\
\code{SpinLockGuard::drop(\&mut self)} &
Libera o lock com ordena\c{c}\~ao \code{Release}. \\
\bottomrule
\end{longtable}

\section{\code{vga.rs} (15 fun\c{c}\~oes)}

\begin{longtable}{@{}p{6.1cm}p{8.4cm}@{}}
\caption{Fun\c{c}\~oes de \code{vga.rs}.}\\
\toprule
\textbf{Assinatura} & \textbf{Descri\c{c}\~ao} \\
\midrule
\endfirsthead
\toprule
\textbf{Assinatura} & \textbf{Descri\c{c}\~ao} \\
\midrule
\endhead
\code{blank\_cell() -> u16} &
C\'elula de espa\c{c}o com atributo padr\~ao; \code{const fn}. \\
\code{Console::new() -> Console} &
Inicializa coluna, linha e cor em contexto \code{const}. \\
\code{Console::clear(\&mut self)} &
Preenche as 2.000 c\'elulas e reposiciona o cursor. \\
\code{Console::write\_byte(\&mut self, byte: u8)} &
Interpreta \code{\textbackslash n}, \code{\textbackslash r},
\code{\textbackslash t}, backspace e ASCII imprim\'ivel. \\
\code{Console::write\_char(\&mut self, ch: char)} &
Grava uma c\'elula e avan\c{c}a a coluna. \\
\code{Console::write\_string(\&mut self, text: \&str)} &
Itera sobre os bytes de uma string. \\
\code{Console::newline(\&mut self)} &
Avan\c{c}a a linha, com rolagem na \'ultima. \\
\code{Console::backspace(\&mut self)} &
Recua uma posi\c{c}\~ao e apaga a c\'elula. \\
\code{Console::scroll(\&mut self)} &
Move o conte\'udo uma linha para cima. \\
\code{Console::update\_cursor(\&self)} &
Escreve a posi\c{c}\~ao na CRTC. \\
\code{<Console as fmt::Write>::write\_str(\&mut self, text: \&str) -> fmt::Result} &
Implementa\c{c}\~ao de \code{fmt::Write} para o console. \\
\code{vga::init()} &
Limpa a tela e configura o cursor de hardware. \\
\code{vga::clear()} &
Limpa a tela a partir do comando \code{clear}. \\
\code{<ConsoleWriter as fmt::Write>::write\_str(\&mut self, text: \&str) -> fmt::Result} &
Escreve no VGA (com lock) e depois na serial. \\
\code{vga::\_print(args: fmt::Arguments<'\_>)} &
Ponto de entrada das macros \code{print!}/\code{println!}. \\
\bottomrule
\end{longtable}

\section{\code{keyboard.rs} (9 fun\c{c}\~oes)}

\begin{longtable}{@{}p{5.3cm}p{9.2cm}@{}}
\caption{Fun\c{c}\~oes de \code{keyboard.rs}.}\\
\toprule
\textbf{Assinatura} & \textbf{Descri\c{c}\~ao} \\
\midrule
\endfirsthead
\toprule
\textbf{Assinatura} & \textbf{Descri\c{c}\~ao} \\
\midrule
\endhead
\code{build\_keymap() -> [u8; 128]} &
Constr\'oi a tabela de scancodes em tempo de compila\c{c}\~ao. \\
\code{apply\_shift(byte: u8) -> u8} &
Aplica Shift a letras, d\'igitos e pontua\c{c}\~ao. \\
\code{wait\_input\_clear() -> bool} &
Espera o buffer de entrada esvaziar. \code{unsafe}. \\
\code{wait\_output\_full() -> bool} &
Espera um byte na sa\'ida. \code{unsafe}. \\
\code{flush\_output()} &
Descarta at\'e 64 bytes pendentes. \code{unsafe}. \\
\code{send\_command(command: u8)} &
Envia comando ao 8042. \code{unsafe}. \\
\code{send\_data(data: u8)} &
Envia dado ao teclado. \code{unsafe}. \\
\code{init()} &
Configura tradu\c{c}\~ao, IRQ, clock e varredura. \\
\code{read\_char() -> Option<u8>} &
L\^e e traduz um scancode; devolve \code{None} se n\~ao h\'a tecla. \\
\bottomrule
\end{longtable}

\section{\code{serial.rs} (4 fun\c{c}\~oes)}

\begin{longtable}{@{}p{4.8cm}p{9.7cm}@{}}
\caption{Fun\c{c}\~oes de \code{serial.rs}.}\\
\toprule
\textbf{Assinatura} & \textbf{Descri\c{c}\~ao} \\
\midrule
\endfirsthead
\toprule
\textbf{Assinatura} & \textbf{Descri\c{c}\~ao} \\
\midrule
\endhead
\code{init()} &
Configura COM1 em 38400 8N1 e detecta presen\c{c}a. \\
\code{read\_byte() -> Option<u8>} &
Leitura n\~ao bloqueante de um byte. \\
\code{write\_byte(byte: u8)} &
Envia um byte, esperando o THR. Interna. \\
\code{write\_str(text: \&str)} &
Envia string com convers\~ao de \code{\textbackslash n} para CRLF. \\
\bottomrule
\end{longtable}

\section{\code{shell.rs} (4 fun\c{c}\~oes)}

\begin{longtable}{@{}p{4.8cm}p{9.7cm}@{}}
\caption{Fun\c{c}\~oes de \code{shell.rs}.}\\
\toprule
\textbf{Assinatura} & \textbf{Descri\c{c}\~ao} \\
\midrule
\endfirsthead
\toprule
\textbf{Assinatura} & \textbf{Descri\c{c}\~ao} \\
\midrule
\endhead
\code{run() -> !} &
La\c{c}o eterno de prompt, leitura e execu\c{c}\~ao. \\
\code{trim(bytes: \&[u8]) -> \&[u8]} &
Remove espa\c{c}os ASCII das extremidades. \\
\code{as\_text(bytes: \&[u8]) -> \&str} &
Converte bytes em \code{\&str} ou texto de erro. \\
\code{execute(line: \&[u8])} &
Compara e executa o comando. \\
\bottomrule
\end{longtable}

\chapter{Mapa de portas e mem\'oria}
\label{ap:portas}

\section{Portas de I/O}

\begin{table}[H]
\centering
\caption{Portas de I/O usadas pelo kernel.}
\label{tab:portmap}
\begin{tabular}{@{}llp{7.5cm}@{}}
\toprule
\textbf{Porta} & \textbf{Dispositivo} & \textbf{Uso} \\
\midrule
\code{0x60} & 8042 dados & Scancodes e dados ao teclado \\
\code{0x64} & 8042 status/comando & Status, comandos, reset \\
\code{0x3D4} & VGA CRTC \'indice & Sele\c{c}\~ao de registrador \\
\code{0x3D5} & VGA CRTC dados & Posi\c{c}\~ao e forma do cursor \\
\code{0x3F8--0x3FF} & UART 16550 COM1 & Dados, LSR, divisor, FIFO \\
\bottomrule
\end{tabular}
\end{table}

\section{Mapa de mem\'oria}

\begin{table}[H]
\centering
\caption{Regi\~oes de mem\'oria relevantes.}
\label{tab:memmap}
\begin{tabularx}{\linewidth}{@{}lX@{}}
\toprule
\textbf{Endere\c{c}o} & \textbf{Uso} \\
\midrule
\code{0x00000000--0x000FFFFF} & Mem\'oria baixa, BIOS, estruturas do GRUB. \\
\code{0x00100000} & In\'icio de \code{.boot} (stub, cabe\c{c}alho, GDT,
tabelas de p\'agina, pilha de 8 KiB). \\
\code{0x00107000} & \emph{Load address} de \code{.text} (virtual alto). \\
\code{0x00109518} & \emph{Load address} de \code{.rodata}. \\
\code{0x0010A000} & \emph{Load address} de \code{.data}. \\
\code{0x0010A020} & \emph{Load address} de \code{.bss}; inclui a pilha de
32 KiB do kernel. \\
\code{0x000B8000} & Framebuffer de texto VGA (80x25, 4.000 bytes). \\
\code{0xFFFFFFFF80000000} & Base virtual do kernel (metade alta). \\
\code{0xFFFFFFFF80107000} & In\'icio virtual de \code{.text}. \\
\bottomrule
\end{tabularx}
\end{table}

\begin{nota}[Endere\c{c}os observados]
Os valores de \code{.text}, \code{.rodata}, \code{.data} e \code{.bss}
correspondem \`a compila\c{c}\~ao de refer\^encia. Eles podem mudar com
otimiza\c{c}\~oes, vers\~ao do compilador ou altera\c{c}\~oes no c\'odigo.
O que n\~ao muda \'e o esquema: stub em 1 MiB, kernel em
\code{0xFFFFFFFF80000000}.
\end{nota}

\chapter{Gloss\'ario}
\label{ap:glossario}

\begin{description}[style=nextline,leftmargin=2em]
  \item[BIOS] Firmware legado de inicializa\c{c}\~ao de PCs, respons\'avel
        por detectar hardware e iniciar o bootloader.
  \item[Bootloader] Programa que carrega o kernel; aqui, o GRUB.
  \item[CR0/CR3/CR4] Registradores de controle da CPU x86; controlam
        pagina\c{c}\~ao, modo protegido e features.
  \item[CRTC] Controlador de tubo de raios cat\'odicos da VGA; controla
        cursor e timing de v\'ideo.
  \item[El Torito] Padr\~ao para tornar CDs/DVDs inicializ\'aveis.
  \item[GDT] Global Descriptor Table; define segmentos de c\'odigo e dados.
  \item[GOP] Graphics Output Protocol, interface de framebuffer em UEFI.
  \item[IDT] Interrupt Descriptor Table; define handlers de interrup\c{c}\~oes
        e exce\c{c}\~oes.
  \item[LME] Long Mode Enable, bit do registrador EFER que habilita o modo
        de 64 bits.
  \item[LSR] Line Status Register da UART; indica dados dispon\'iveis e
        transmissor vazio.
  \item[Multiboot2] Especifica\c{c}\~ao de protocolo entre bootloader e
        kernel; usada pelo GRUB.
  \item[PML4/PDPT/PD] N\'iveis da tabela de p\'aginas x86\_64.
  \item[Polling] Consulta repetida de hardware, sem interrup\c{c}\~oes.
  \item[Ring 0] N\'ivel de privil\'egio mais alto da CPU; c\'odigo de kernel.
  \item[Scancode] C\'odigo enviado pelo teclado para representar uma tecla.
  \item[Spinlock] Lock que faz a CPU girar at\'e a regi\~ao cr\'itica ser
        liberada.
  \item[Stub] C\'odigo inicial em assembly que prepara a CPU antes do Rust.
  \item[Triple fault] Falha de exce\c{c}\~ao durante o tratamento de outra,
        que faz a CPU resetar.
  \item[UART] Universal Asynchronous Receiver/Transmitter; controlador
        serial.
  \item[VGA texto] Modo de v\'ideo legado com matriz de caracteres em
        \code{0xB8000}.
\end{description}

\begin{thebibliography}{9}

\bibitem{intel}
Intel Corporation.
\newblock \emph{Intel 64 and IA-32 Architectures Software Developer's
Manual}.
\newblock Volume 3A: System Programming Guide, cap\'itulos sobre modo
protegido, pagina\c{c}\~ao e long mode.

\bibitem{multiboot2}
Free Software Foundation.
\newblock \emph{Multiboot2 Specification}.
\newblock Dispon\'ivel em \url{https://www.gnu.org/software/grub/manual/multiboot2/}.

\bibitem{linux}
Linux Kernel Organization.
\newblock \emph{Linux Kernel 7.3.1 Source Tree}.
\newblock Arquivos \code{arch/x86}, \code{drivers/video/console/vgacon.c},
\code{drivers/input}, \code{drivers/tty/serial/8250} e \code{kernel/}.

\bibitem{osdev}
OSDev Wiki.
\newblock \emph{PS/2 Keyboard}, \emph{Serial Ports}, \emph{VGA Text Mode} e
\emph{Setting Up Long Mode}.
\newblock Dispon\'ivel em \url{https://wiki.osdev.org/}.

\bibitem{rust}
The Rust Project.
\newblock \emph{The Rust Programming Language} e \emph{Rust Platform Support:
x86\_64-unknown-none}.
\newblock Dispon\'ivel em \url{https://doc.rust-lang.org/}.

\bibitem{qemu}
QEMU Project.
\newblock \emph{QEMU Emulator User Documentation}.
\newblock Dispon\'ivel em \url{https://www.qemu.org/documentation/}.

\bibitem{grub}
GNU GRUB.
\newblock \emph{GNU GRUB Manual}.
\newblock Dispon\'ivel em \url{https://www.gnu.org/software/grub/manual/grub/}.

\end{thebibliography}
